% ==========================================
% LatexRefinement Step Output: step4-03-03-25-tuesday-all-offset-final.tex
% Model: gemini-3-flash-preview
% Temperature: 0.33
% TopP: 0.95
% TopK: 40
% MaxOutputTokens: 65535
% ThinkingBudget: 32765
% ThinkingLevel: HIGH
% Prompt Tokens: 36'911
% Candidates Tokens: 12'903 (inkl. Thinking Tokens)
% Cached Tokens: 32'746
% Processed on: 2026-07-15 11:48:39
% ==========================================

\lecturechapter[3]{Dienstag}{03. Mär}{3. März 2026}{Peano-Arithmetik und formale Beweise}

\begin{nice-box}[Vorlesungskontext und Übersicht]
In dieser Vorlesung setzen wir unsere Reise durch die mathematische Logik fort. Nachdem wir in der letzten Woche die Grundlagen formaler Beweise und wichtige Metatheoreme wie das Deduktionstheorem behandelt haben, widmen wir uns heute zwei zentralen Themen:
\begin{itemize}
    \item \textbf{Peano-Arithmetik (PA):} Wir untersuchen die Axiomatisierung der natürlichen Zahlen, führen das Induktionsschema ein und vollziehen erste (semi-)formale Beweise innerhalb dieses Systems.
    \item \textbf{Modelltheorie:} Wir vollziehen den Übergang von der rein syntaktischen Ebene (Zeichenketten und Umformungsregeln) zur semantischen Ebene (Bedeutung, Wahrheit und Strukturen). Dabei führen wir den Begriff der $\mathcal{L}$-Struktur, der Interpretation und die Tarski-Wahrheitsdefinition ein.
\end{itemize}
\end{nice-box}

\section{Rückblick und Einordnung}

\begin{spoken-clean}[00:00:00 - 00:02:15]
Hallo zusammen und herzlich willkommen zu einer weiteren Woche Grundstrukturen. Wir gehen weiter auf unserer kleinen Exkursion in die mathematische Logik. Was haben wir letzte Woche gesehen?

Wir haben zuerst den Modus Ponens (\textbf{MP}) und die Verallgemeinerung (\textbf{V}) gesehen. Das sind zwei Methoden, um aus bestehenden Formeln neue Formeln abzuleiten. Damit haben wir dann definiert, was ein formaler Beweis ist, und haben ein paar formale Beweise geführt. Das machen Sie auch diese Woche --- oder haben es hoffentlich schon gemacht --- auf den Übungsblättern, um noch ein bisschen zu üben. Ich glaube, es ist einfach etwas, was man einmal ein klein wenig machen muss --- nicht allzu viel. Es kann auch ganz unterhaltsam sein, da etwas herumzuknobeln, aber auf die Dauer ist es auch mühsam.

Dann haben wir noch diese Metatheoreme gesehen, die manchmal hilfreich sind, um formale Beweise zu führen: das Deduktionstheorem und der Satz über logische Äquivalenz. Und am Schluss haben wir noch Beispiele von Axiomensystemen angeschaut.
\end{spoken-clean}

\subsection{Rückblick: Letzte Woche}
\begin{math-stroke}
\textbf{Gesehen:}
\begin{itemize}
    \item \textbf{(MP)} und \textbf{(V):} Modus Ponens und Verallgemeinerung.
    \item \textbf{Formale Beweise:} Definition und Durchführung.
    \item \textbf{Deduktionstheorem \& Satz über logische Äquivalenz:} Wichtige Metatheoreme zur Vereinfachung von Beweisen.
    \item \textbf{Beispiele von Axiomensystemen:} Verschiedene mathematische Theorien in prädikatenlogischer Form.
\end{itemize}
\end{math-stroke}

\section{Dichte lineare Ordnungen}

\begin{spoken-clean}[00:02:15 - 00:04:20]
Genau, da möchte ich noch kurz eines erwähnen, das wir noch nicht besprochen haben: die dichten linearen Ordnungen. Das sind die folgenden fünf Formeln, die unsere Axiome bilden.

Das erste Axiom, \textbf{DLO$_0$}, sagt: Für alle $x$ gilt nicht $x < x$. Wir haben einfach ein einzelnes zweistelliges Relationssymbol, das ist einfach \qt{kleiner}. Wir wollen eine totale Ordnung haben; also wollen wir, dass $x$ nie kleiner ist als $x$ selbst. Dann wollen wir, dass für alle $x, y, z$ gilt: Wenn $x < y$ und $y < z$, dann ist auch $x < z$. Das ist genau das, was wir uns vorstellen, wenn wir diese Formel hinschreiben --- dass die Relation transitiv ist ---, aber hier ist es eben erst einmal nur eine Formel. Dann wollen wir, dass für alle $x$ und $y$ entweder $x < y$ oder $y < x$ oder $x = y$ gilt. Es ist natürlich, dass es keine anderen Möglichkeiten gibt: eine totale Ordnung.

Dann wollen wir, dass die Ordnung dicht ist. Das heißt, wir wollen, dass zwischen zwei Elementen quasi immer ein drittes liegt. Also: Für alle $x, y$ existiert immer ein $z$, so dass falls $x < y$ gilt, dann auch $x < z$ und $z < y$ gilt. Das ist diese Dichtheit, wie zum Beispiel bei den reellen Zahlen: Wenn Sie zwei reelle Zahlen haben, dann haben Sie immer noch ein anderes Element dazwischen. Und das letzte Axiom besagt, dass es keine größten und kleinsten Elemente gibt. Also: Für alle $x$ existiert ein $y$ und es existiert ein $z$, so dass $y < x$ und $x < z$ gilt.

Okay, das ist einfach ein wichtiger Begriff, den wir auch noch in den Übungen antreffen werden. Es ist ein weiteres Beispiel für eine Menge von nicht-logischen Axiomen, die man anwenden kann, um entsprechende Theorien zu studieren.
\end{spoken-clean}

\subsection{Dichte lineare Ordnungen (DLO)}
\begin{math-stroke}
Die Signatur der Theorie der dichten linearen Ordnungen ist $\mathcal{L}_{\text{DLO}} := \{ < \}$, wobei $<$ ein $2$-stelliges Relationssymbol ist.

Die Axiome der Theorie der dichten linearen Ordnungen sind:
\begin{align}
    \mathbf{DLO_0}\colon &\quad \forall x \, \neg(x < x) && (\text{irreflexiv}) \nonumber \\
    \mathbf{DLO_1}\colon &\quad \forall x \forall y \forall z \, (x < y \land y < z \to x < z) && (\text{transitiv}) \nonumber \\
    \mathbf{DLO_2}\colon &\quad \forall x \forall y \, (x < y \lor x = y \lor y < x) && (\text{linear/total}) \nonumber \\
    \mathbf{DLO_3}\colon &\quad \forall x \forall y \, (x < y \to \exists z \, (x < z \land z < y)) && (\text{dicht}) \nonumber \\
    \mathbf{DLO_4}\colon &\quad \forall x \exists y \exists z \, (y < x \land x < z) && (\text{randlos}) \nonumber
\end{align}
\end{math-stroke}

\section{Peano-Arithmetik}

\begin{meta-note}[Tafelübergang]
Der Dozent schaltet den Projektor aus und wechselt zur Tafel, um die Peano-Arithmetik einzuführen.
\end{meta-note}

\begin{spoken-clean}[00:04:20 - 00:06:29]
Gut, dann schauen wir uns noch ein weiteres System an. Das möchte ich noch ein bisschen ausführen, deswegen machen wir das an der Wandtafel: die Peano-Arithmetik. Sie ist benannt nach Giuseppe Peano, einem italienischen Mathematiker, der in der zweiten Hälfte des 19. Jahrhunderts gelebt hat. Er hat sehr viel Einflussreiches in der Logik geleistet. Ich glaube, viele bekannte Mengennotationen gehen auf ihn zurück --- zum Beispiel, dass man diesen Haken für den Durchschnitt von Mengen schreibt und das entsprechende Symbol für die Vereinigung. Er hat auch 1897 auf dem Internationalen Mathematikerkongress hier in Zürich einen Vortrag gehalten.

Er hat diese Peano-Axiome formuliert. Die Idee war eigentlich, die natürlichen Zahlen in Axiome zu fassen. So wie wir sie formulieren, sind es zwar nicht exakt die ganzen natürlichen Zahlen, aber es geht in die Richtung --- Sie werden es sehen. Haben Sie die Peano-Axiome schon einmal in einer anderen Vorlesung gesehen? Manchmal macht man das in der Linearen Algebra oder in der Analysis. Nein? Gut, dann machen wir es hier.
\end{spoken-clean}

\subsection{Signatur der Peano-Arithmetik}

\begin{spoken-clean}[00:06:29 - 00:07:37]
Hier haben wir die Signatur der Peano-Arithmetik: $0, S, +$ und $\cdot$. Dabei ist $0$ ein Konstantensymbol. Dann ist $S$ ein einstelliges Funktionssymbol. Wir haben $+$, ein zweistelliges Funktionssymbol, und $\cdot$ ist ebenfalls ein zweistelliges Funktionssymbol. Wenig erstaunlich, das ist einfach das, was wir bereits kennen.

Das Einzige, was vielleicht ein bisschen ungewohnt ist, ist dieses $S$ hier. Der Name von $S$ ist Nachfolgerfunktion. Die Intuition dahinter ist, die natürlichen Zahlen zu konstruieren, indem man quasi zählt. Man beginnt mit $0$, das ist einfach eine Konstante, von der man weiß, dass es sie gibt. Und dann gibt es eben $S$. Die Idee dahinter ist: Man macht quasi $+1$. Man beginnt also bei $0$, dann macht man $+1$. Wenn man nochmals $S$ anwendet, kommt man zu $+2$ und so weiter. Je nachdem, wie oft man $S$ anwendet, gelangt man zu einer bestimmten Zahl.
\end{spoken-clean}

\begin{math-stroke}[Signatur der Peano-Arithmetik]
Die Signatur der Peano-Arithmetik ist:
\[
\mathcal{L}_{\text{PA}} := \{0, S, +, \cdot\}
\]
wobei:
\begin{itemize}
    \item \textbf{$0$} ein Konstantensymbol ist.
    \item \textbf{Funktion:} \newterm{Nachfolgerfunktion} $S$ ist ein $1$-stelliges Funktionssymbol.
    \item \textbf{$+$} ein $2$-stelliges Funktionssymbol ist.
    \item \textbf{$\cdot$} ein $2$-stelliges Funktionssymbol ist.
\end{itemize}
\end{math-stroke}

\subsection{Die Axiome der Peano-Arithmetik}

\begin{spoken-clean}[00:07:37 - 00:11:00]
Die Axiome der Peano-Arithmetik sind nun die folgenden Formeln. Wir haben $\mathbf{PA_0}$. Das besagt: Es existiert kein $x$, so dass $S(x) = 0$ ist. Die Idee dahinter --- auch wenn das hier nur Formeln sind --- ist: $0$ ist nicht der Nachfolger von irgendeinem anderen Element.

Gut, dann haben wir $\mathbf{PA_1}$. Das sagt uns, dass diese Nachfolgerfunktion injektiv sein soll. Wie kann man das formulieren? Für alle $x$ und für alle $y$ gilt: $S(x) = S(y)$ impliziert $x = y$. Das ist genau die Formel, mit der man ausdrücken möchte, dass $S$ injektiv ist.

Dann haben wir $\mathbf{PA_2}$. Es gibt insgesamt sieben von diesen Axiomen. Das besagt, dass für alle $x$ gilt: $x + 0 = x$. Jetzt regeln wir die Addition. Wir wollen, dass $0$ das rechtsneutrale Element ist. Dann haben wir $\mathbf{PA_3}$. Das sagt: Für alle $x$ und für alle $y$ haben wir $x$ plus den Nachfolger von $y$. Das kann man quasi als $x + (y + 1)$ lesen, und das soll dasselbe sein wie der Nachfolger von $x + y$. Also: Ob man das $S$ einfach auf das hintere Element anwendet oder auf beide zusammen, es kommt dasselbe heraus.

Dann haben wir $\mathbf{PA_4}$. Jetzt kommt noch die Multiplikation. Das sagt: Für alle $x$ gilt, dass $x \cdot 0$ --- wir schreiben den Punkt natürlich dazwischen --- gleich $0$ ist. Dann kommt noch $\mathbf{PA_5}$. Das sagt: Für alle $x$ und für alle $y$ ist $x$ mal der Nachfolger von $y$ dasselbe wie $(x \cdot y) + x$. Gut. Soweit also diese ersten fünf Axiome.
\end{spoken-clean}

\begin{math-stroke}[Die Axiome der Peano-Arithmetik]
\begin{axioms}[Peano-Arithmetik]
Die nachstehenden Axiome bilden den Kern der Peano-Arithmetik:
\begin{align}
\mathbf{PA_0}\colon &\quad \neg \exists x \, (S(x) = 0) \label{eq:pa0}\\
\mathbf{PA_1}\colon &\quad \forall x \forall y \, (S(x) = S(y) \to x = y) \label{eq:pa1}\\
\mathbf{PA_2}\colon &\quad \forall x \, (x + 0 = x) \label{eq:pa2}\\
\mathbf{PA_3}\colon &\quad \forall x \forall y \, (x + S(y) = S(x + y)) \label{eq:pa3}\\
\mathbf{PA_4}\colon &\quad \forall x \, (x \cdot 0 = 0) \label{eq:pa4}\\
\mathbf{PA_5}\colon &\quad \forall x \forall y \, (x \cdot S(y) = (x \cdot y) + x) \label{eq:pa5}
\end{align}
\end{axioms}
\end{math-stroke}

\subsection{Das Induktionsaxiom}

\begin{spoken-clean}[00:11:00 - 00:13:28]
Und jetzt kommt noch ein Axiomenschema. Was war ein Axiomenschema noch gleich? Das bedeutet, es ist eine riesig große Menge von Axiomen; für jede Formel haben wir ein Axiom. Und das ist dieses Induktionsaxiom. Wir sagen: Wenn $\varphi$ eine Formel mit dieser Signatur ist und $\nu$ eine freie Variable in dieser Formel, dann haben wir $\mathbf{PA_6}$.

Das ist eben dieses ganze Schema. Es besagt im Prinzip: Wenn eine Formel für $0$ stimmt und wenn sie für $x$ stimmt, dann stimmt sie auch für den Nachfolger --- und damit für alle. Wir wollen also einfach das Prinzip des Induktionsbeweises in ein Axiomenschema stecken.

Das heißt: Wenn $\varphi(0)$ gilt und für alle $\nu$ gilt, dass $\varphi(\nu)$ die Formel $\varphi(S(\nu))$ impliziert, dann folgt aus diesen zwei Dingen, dass $\varphi(\nu)$ für alle $\nu$ gilt. Das ist genau das, was uns ein Induktionsbeweis sagt. Ich schreibe einfach hin: Das ist das Induktionsaxiom.
\end{spoken-clean}

\begin{math-stroke}[Induktionsschema]
\textbf{Induktionsschema:}
Sei $\varphi$ eine $\mathcal{L}_{\text{PA}}$-Formel und $\nu$ eine freie Variable in $\varphi$.
\begin{equation}
\label{eq:pa6}
\mathbf{PA_6}\colon \quad \bigl( \varphi(0) \land \forall \nu \, (\varphi(\nu) \to \varphi(S(\nu))) \bigr) \to \forall \nu \, \varphi(\nu)
\end{equation}

\begin{explanation-of-steps}[Warum ein Schema? (und Notationshinweis)]
In der Prädikatenlogik erster Stufe können wir nicht über Mengen von Formeln quantifizieren. Wir können also nicht ein einzelnes Axiom schreiben, das besagt: \qt{Für alle Formeln $\varphi$ gilt...}. Stattdessen führen wir für jede mögliche Formel $\varphi$ ein eigenes Axiom ein, das die Struktur von \eqref{eq:pa6} besitzt. In der Logik zweiter Ordnung könnte man dies mit einem einzigen Axiom regeln, was jedoch andere metallogische Probleme aufwirft.

(Zudem ist die Schreibweise $\varphi(0)$ und $\varphi(S(\nu))$ eine informelle Kurznotation für die formale Substitution $\varphi(\nu/0)$ und $\varphi(\nu/S(\nu))$.)
\end{explanation-of-steps}
\end{math-stroke}

\section{Formale und semi-formale Beweise}

\begin{spoken-clean}[00:13:28 - 00:14:48]
In gewissen Vorlesungen --- oft wird es in der Linearen Algebra oder in der Analysis diskutiert --- finden Sie auch leicht andere Versionen. Ich glaube, vor allem die Art und Weise, wie ich das Induktionsaxiom geschrieben habe, steht wirklich im Kontext dieser Logik erster Ordnung. Das heißt, wir müssen wirklich für jede Formel ein Axiom einführen, weil wir nicht über die Menge der Formeln quantifizieren können. Wir können also nicht ein Axiom schreiben: \qt{Für alle Formeln $\varphi$ gilt das}. Sondern wir müssen jede einzelne Formel dort einsetzen.

In der Logik zweiter Ordnung darf man auch über Formeln quantifizieren; dann kann man das mit einem einzigen Axiom regeln. Aber Logik zweiter Ordnung hat dann wiederum andere Probleme, deswegen machen wir das hier so. Wie wir später noch erwähnen werden, definiert uns das in der Logik erster Ordnung nicht eindeutig die natürlichen Zahlen. Man kann die natürlichen Zahlen gar nicht in Logik erster Ordnung charakterisieren, wie wir sehen werden. Aber es ist auch kein Problem, dass man das nicht kann.
\end{spoken-clean}

\begin{didactic-insight}[Die Grenzen der Logik erster Stufe]
Ein tiefes Resultat der mathematischen Logik ist, dass die Peano-Axiome in der Logik erster Stufe (First-Order Logic) die Struktur der natürlichen Zahlen nicht eindeutig festlegen. Es existieren sogenannte \newterm{Nicht-Standard-Modelle} der Arithmetik, die neben den üblichen natürlichen Zahlen noch weitere, \qt{unendlich große} Elemente enthalten. Dennoch reicht die Peano-Arithmetik aus, um fast die gesamte klassische Zahlentheorie formal zu entwickeln.
\end{didactic-insight}

\begin{spoken-clean}[00:14:48 - 00:17:30]
Gut, das ist wichtig. Deswegen gibt es diese Woche auf dem Übungsblatt nochmals zwei Übungen, in denen Sie mithilfe dieser Axiome einen formalen Beweis ausführen sollen. Da kommen so ganz einfache Sachen vor. Zum Beispiel eine Übung --- es ist nicht schlecht, das einmal wirklich formal sauber zu machen ---: Man soll zeigen, dass aus der Peano-Arithmetik folgt, dass der Nachfolger von $0$ plus der Nachfolger von $0$ dasselbe ist wie der zweifache Nachfolger von $0$. In anderen Worten: $1 + 1 = 2$.

Vielleicht ist es nicht schlecht, das im Laufe des Mathestudiums einmal wirklich im Detail zu beweisen. Das, und dann gibt es noch etwas Ähnliches zu beweisen. Danach hören wir aber auch auf mit diesen formalen Beweisen, weil das eben mühsam ist.

Ich habe noch eine Bemerkung und einen Link auf die Moodle-Seite gestellt. Die Frage ist ja: Macht man dieses ganze formale Beweisen überhaupt? Wenn man heutzutage Mathematik betreibt, macht eigentlich niemand formale Beweise. Fast niemand führt die Beweise auf die Axiome zurück, weil das jeglichen Rahmen einer Publikation sprengen würde. Schlussendlich sind es auch nur die Logiker, die im Alltag wirklich so nachdenken. Aber es ist natürlich trotzdem spannend, gewisse Beweise wirklich auf ein solides Fundament zu stellen. Mit dem Computer kann man das ja machen. Es gibt tatsächlich eine Community, die fleißig daran arbeitet, große Teile der Mathematik auf solide Fundamente zu stellen.
\end{spoken-clean}

\subsection{Formale Verifikation und LEAN}
\begin{math-stroke}
\begin{center}
\begin{color-box}{blue}[LEAN Theorem Prover]
Eines der am weitesten verbreiteten Computerprogramme zur formalen Verifikation ist \newterm{LEAN}. Es verwendet eine etwas andere Logik (Typentheorie), die jedoch logisch im Wesentlichen äquivalent zur hier eingeführten Prädikatenlogik erster Stufe ist.
\end{color-box}
\end{center}

\begin{explanation-of-steps}[Wie funktioniert Computer-gestütztes Beweisen?]
Das System besteht meist aus zwei Teilen:
\begin{enumerate}[label=\textbf{(\arabic*)}]
    \item Einem Interface, das die Eingabe von mathematischen Aussagen in Code-Form ermöglicht.
    \item Einem \qt{Kernel}, der den Beweis auf die grundlegendsten logischen Schritte reduziert und kontrolliert.
\end{enumerate}
Es gibt große Forschungsprojekte, die versuchen, monumentale Sätze wie den Großen Satz von Fermat vollständig in LEAN zu formalisieren.
\end{explanation-of-steps}
\end{math-stroke}

\begin{spoken-clean}[00:17:30 - 00:20:50]
Jedenfalls, wenn Sie das interessiert und Sie ein bisschen Spaß an diesen formalen Beweisen haben, würde ich empfehlen: Es gibt eine einfache Einführung mit viel \qt{Gamification}, das \qt{Natural Numbers Game}. Ich habe einen Link dazu geteilt. Dort gibt es ein kurzes Tutorial zu den grundlegenden LEAN-Begriffen, und das Spiel führt Sie dann dadurch, dass Sie viele kleine Eigenschaften über natürliche Zahlen beweisen. Es beginnt auch damit zu beweisen, dass $1 + 1 = 2$ ist und so weiter. Falls Sie Spaß daran haben, dürfen Sie das gerne freiwillig weitermachen. Es ist selbstverständlich weit davon entfernt, prüfungsrelevant zu sein.

Gut. Was wir jetzt noch kurz erwähnen wollen, sind \newterm{semi-formale Beweise}. Das ist keine formale Definition, wie man erwarten könnte. Die Frage ist immer: Wenn man einen Beweis führt, wie präzise ist man? Die Idee bei semi-formalen Beweisen ist einfach --- und das werden wir in Zukunft auch tun, wenn wir über Logik reden ---, dass man die Schritte weglässt, in denen man die logischen Axiome anwendet. Man könnte sagen: Die logischen Axiome sind so grundlegend, dass wir sie im Prinzip ausführen könnten, wenn wir uns nur lange genug hinsetzen würden. Deswegen lassen wir sie weg, damit das Ganze übersichtlicher wird. Das ist dann ein semi-formaler Beweis. Man verwendet nur noch die Axiome der Theorie, mit der wir arbeiten.

Es ist immer noch mühsam genug, aber eben. In einem semi-formalen Beweis lassen wir die Schritte aus, in denen wir logische Axiome anwenden, und verwenden stattdessen oft natürliche Sprache. Das ist natürlich der erste Schritt auf dem Weg zum Schlendrian, aber man kann nicht alles formal machen --- sonst wären Sie heute noch nicht einmal in der zweiten Woche Ihrer Linearen Algebra. Deswegen geht man solche Schritte. Aber es ist nicht so, dass man deswegen jetzt keine Logik mehr betreibt. Man führt nur gewisse Dinge nicht mehr explizit aus. Wenn man natürliche Sprache verwendet, kann man im Hinterkopf behalten, was man eigentlich gerade tut oder tun sollte.\end{spoken-clean}

\subsection{Semi-formale Beweise}
\begin{math-stroke}
\begin{definition}[Semi-formaler Beweis]
Ein \newterm{semi-formaler Beweis} ist eine Beweisführung, bei der:
\begin{itemize}
    \item Die rein logischen Axiom-Anwendungen (wie Instanzen von \textbf{L1}--\textbf{L16}) meist unterdrückt werden.
    \item Nur die theoriespezifischen Axiome (z.\,B. $\mathbf{PA_0}$--$\mathbf{PA_6}$) explizit genannt werden.
    \item Natürliche Sprache zur Strukturierung verwendet wird, wobei jeder Schritt im Prinzip auf die formalen Axiome zurückführbar bleibt.
\end{itemize}
\end{definition}
\end{math-stroke}

\section{Beispielbeweise in der Peano-Arithmetik}

\begin{spoken-clean}[00:20:50 - 00:24:50]
Schauen wir uns einfach einmal ein Beispiel an. Als erstes Beispiel nehmen wir diese Übung hier. Sie sollen das zwar noch formal beweisen --- und Sie werden sehen, das ist mühsam ---, aber wir machen das jetzt einmal semi-formal. Wir zeigen, dass in der Peano-Arithmetik $S(0) + S(0) = S(S(0))$ gilt.

Ein semi-formaler Beweis hier wäre zum Beispiel: Wir haben $S(0) + S(0)$. Jetzt wenden wir \textbf{PA$_3$} an. \textbf{PA$_3$} sagt uns, dass $S(0) + S(0)$ dasselbe ist wie $S(S(0) + 0)$. Das ist einfach gemäß Peano-Axiom 3. Mit \textbf{PA$_2$} folgt nun, dass das wiederum dasselbe ist wie $S(S(0))$, weil $S(0) + 0$ wieder $S(0)$ ergibt. Das ist der semi-formale Beweis und natürlich wesentlich einfacher. Es ist immer noch nicht ganz trivial zu zeigen, dass $1 + 1 = 2$ ist, aber schon deutlich kürzer als direkt aus den Axiomen.

Gut. Noch ein Beispiel, das etwas mühsamer ist: Wir wollen als Nächstes zeigen, dass $S(S(0)) \cdot x = x + x$ gilt. Das ist etwas schwieriger und wir wollen es mit Induktion machen. Hier vielleicht noch zwei Resultate, die Sie in den Übungen ebenfalls semi-formal zeigen: Erstens, dass die Addition assoziativ ist --- also für alle $x, y, z$ gilt $(x + y) + z = x + (y + z)$. Und zweitens, dass die Addition kommutativ ist, also für alle $x, y$ gilt $x + y = y + x$.
\end{spoken-clean}

\subsection{Übung: Addition und Multiplikation}
\begin{math-stroke}
\begin{example}[Addition von Einsen]
$\text{PA} \vdash S(0) + S(0) = S(S(0))$
\end{example}

\begin{short-proof}
\begin{align*}
    S(0) + S(0) &= S(S(0) + 0) && (\text{nach } \mathbf{PA_3}) \\
    &= S(S(0)) && (\text{nach } \mathbf{PA_2})
\end{align*}
\end{short-proof}

\begin{remark}[Assoziativität und Kommutativität]
Auf dem Übungsblatt werden folgende fundamentale Eigenschaften der Addition semi-formal bewiesen:
\begin{itemize}
    \item \textbf{Assoziativität:} $\forall x \forall y \forall z \, ((x + y) + z = x + (y + z))$
    \item \textbf{Kommutativität:} $\forall x \forall y \, (x + y = y + x)$
\end{itemize}
\end{remark}
\end{math-stroke}

\subsection{Beweis der Identität \texorpdfstring{$2 \cdot x = x + x$}{2x = x+x}}

\begin{spoken-clean}[00:24:50 - 00:28:30]
Jetzt wollen wir Induktion anwenden. Vielleicht mache ich das auf dieser Tafel. Ist es okay, wenn ich hier auswische? \inlinemetanote{wartet auf Zustimmung der Studierenden} Gut, machen wir das, damit wir die Axiome dort noch stehen haben. \inlinemetanote{wischt die Tafel}

Wir wollen zeigen: In der Peano-Arithmetik gilt für alle $x$, dass $S(S(0)) \cdot x = x + x$ ist. Wir führen wieder einen semi-formalen Beweis. Dafür definieren wir die Formel $\varphi(x)$, die wir beweisen wollen: $S(S(0)) \cdot x = x + x$. Wir verwenden Induktion. Das heißt, wir zeigen zuerst, dass die Aussage für $x = 0$ stimmt. Dann zeigen wir: Wenn sie für $x$ stimmt, dann stimmt sie auch für den Nachfolger $S(x)$. Dann verwenden wir Axiom \textbf{PA$_6$} und wissen, dass es für alle $x$ gilt.

Zuerst müssen wir beweisen, dass in der Peano-Arithmetik $\varphi(0)$ gilt. Eigentlich müsste man hier $\varphi[0/\nu]$ schreiben, weil wir $x$ durch $0$ substituieren, aber wenn aus dem Kontext klar ist, welche Variable wir substituieren, müssen wir das nicht explizit hinschreiben. Das ist relativ einfach: Wir haben $S(S(0)) \cdot 0$. Das ist $0$ gemäß \textbf{PA$_4$}. Wenn man etwas mit $0$ multipliziert, ergibt das wieder $0$. Dann verwenden wir \textbf{PA$_2$}, was uns sagt, dass das dasselbe ist wie $0 + 0$. Okay, das ist genau das, was wir zeigen wollten. Die Induktionsverankerung stimmt also.
\end{spoken-clean}

\subsubsection{Satz: Multiplikation mit 2}
\begin{math-stroke}
\begin{proposition}
    \label{prop:2x_is_xx}
    $\text{PA} \vdash S(S(0)) \cdot x = x + x$
\end{proposition}

\begin{short-proof}[Beweis durch Induktion]
Wir definieren die Formel $\varphi(x) \equiv S(S(0)) \cdot x = x + x$.

\textbf{Induktionsanfang ($\varphi(0)$):}
Zu zeigen ist $S(S(0)) \cdot 0 = 0 + 0$.
\begin{itemize}
    \item Linke Seite: $S(S(0)) \cdot 0 = 0$ (nach Axiom \ref{eq:pa4}).
    \item Rechte Seite: $0 + 0 = 0$ (nach Axiom \ref{eq:pa2}).
\end{itemize}
Also gilt $\varphi(0)$.
\end{short-proof}
\end{math-stroke}

\begin{spoken-clean}[00:28:30 - 00:31:45]
Gut, und jetzt wollen wir noch den Induktionsschritt zeigen. Wir zeigen: In der Peano-Arithmetik gilt, dass $\varphi(x)$ die Formel $\varphi(S(x))$ impliziert. Wir nehmen an, dass $\varphi(x)$ gilt. Semi-formaler Beweis: Wir nehmen an, dass $\varphi(x)$ gilt, also $S(S(0)) \cdot x = x + x$.

Jetzt formen wir das um. Wir betrachten $S(S(0)) \cdot S(x)$. Gemäß \textbf{PA$_5$} ist das dasselbe wie $(S(S(0)) \cdot x) + S(S(0))$. Da wissen wir nun gemäß Induktionsannahme, dass das dasselbe ist wie $(x + x) + S(S(0))$. Jetzt können wir die Assoziativität und Kommutativität verwenden, die Sie in den Übungen zeigen. Wir können das umschreiben als $(x + S(S(0))) + x$. Und $x + S(S(0))$ ist nach \textbf{PA$_3$} dasselbe wie $S(x + S(0))$, was wiederum $S(S(x + 0))$ ist, also $S(S(x))$. Also haben wir $S(S(x)) + x$. Und das ist wiederum $S(x) + S(x)$, wenn man die Axiome richtig anwendet.

Das zeigt uns also den Induktionsschritt. Zusammen mit dem Induktionsanfang folgt nach dem Induktionsaxiom \textbf{PA$_6$}, dass die Formel für alle $x$ gilt. Gut, das ist wichtig. Deswegen gibt es diese Woche auf dem Übungsblatt nochmals zwei Übungen, in denen Sie wirklich nochmals einen formalen Beweis ausführen sollen.
\end{spoken-clean}

\subsubsection{Induktionsschritt}
\begin{math-stroke}
\begin{short-proof}[Fortsetzung des Beweises]
\textbf{Induktionsschritt ($\varphi(x) \to \varphi(S(x))$):}
Annahme (IV): $S(S(0)) \cdot x = x + x$.
Zu zeigen: $S(S(0)) \cdot S(x) = S(x) + S(x)$.

Es gilt:
\begin{align*}
    S(S(0)) \cdot S(x) &= (S(S(0)) \cdot x) + S(S(0)) && (\text{nach } \mathbf{PA_5}) \\
    &= (x + x) + S(S(0)) && (\text{nach IV}) \\
    &= (x + x) + S(S(0)) && (\text{Assoziativität/Kommutativität}) \\
    &= S(x) + S(x) && (\text{nach } \mathbf{PA_3} \text{ und } \mathbf{PA_2})
\end{align*}
Damit ist der Induktionsschritt gezeigt. Nach dem Induktionsschema \eqref{eq:pa6} folgt die Behauptung für alle $x$.
\end{short-proof}
\end{math-stroke}

\begin{ai-note}[{Lücke im Induktionsschritt --- Sonnet 5.x / Medium}]
Die dritte Zeile der obigen `align*`-Kette ist textidentisch mit der zweiten (kein tatsächlicher Umformungsschritt), und die Endbegründung "nach $\mathbf{PA_3}$ und $\mathbf{PA_2}$" überspringt die eigentliche Umgruppierung. Vollständig: $(x+x)+S(S(0)) = (x+S(S(0)))+x$ (Kommutativität/Assoziativität) $= S(S(x))+x$ (zweimal $\mathbf{PA_3}$ und $\mathbf{PA_2}$) $= S(x)+S(x)$ (Kommutativität und $\mathbf{PA_3}$). Das Endresultat bleibt korrekt.
\end{ai-note}

\section{Modelltheorie}

\begin{spoken-clean}[00:31:45 - 00:33:15]
Gut. So viel zur Peano-Arithmetik für den Moment und auch zu semi-formalen Beweisen. Gibt es dazu noch Fragen?

Sonst gehen wir jetzt zum nächsten Kapitel. Da kommen die Modelle --- das ist der Bereich der Modelltheorie. Ein neues Kapitel. \inlinemetanote{wischt die Tafel} Vielleicht muss man philosophisch ein bisschen aufpassen, was man genau macht. Was wir bis jetzt getan haben, ist: Wir haben rein sprachlich arbeitet. Wir haben diese ganze Syntax aufgebaut und ein Alphabet festgelegt, wie man all die Sachen umformen darf. Aber es war eine reine Sprache, die wir entwickelt haben. Sie hat a priori keine Bedeutung, sondern folgt einfach syntaktischen Regeln.
\end{spoken-clean}

\begin{spoken-clean}[00:33:15 - 00:35:40]
Die Idee ist: Wir wollen dem Ganzen nun Sinn geben. Wir wollen von dieser Syntaktik auf die semantische Ebene wechseln. Auf der semantischen Ebene schauen wir uns tatsächlich nicht einfach nur die Gruppenaxiome an, sondern tatsächlich Gruppen. Wir wollen dem Ganzen Bedeutung verleihen und über \qt{wahr} und \qt{falsch} sprechen können.

Aber das bedeutet natürlich auch wieder, dass man einen Schritt zurücktritt. Wir können das nicht in derselben Sprache machen, in der wir die Syntax aufgebaut haben, sonst würde sich die Schlange irgendwann in den Schwanz beißen. Das heißt, wir werden jetzt quasi eine naive Mengenlehre verwenden. Wir werden von Mengen sprechen, von Teilmengen und von Funktionen, aber ohne dabei allzu komplizierte Mengenbegriffe einzuführen.

Es bleibt natürlich die Frage, wo genau das Ganze \qt{lebt}. Später behandeln wir noch die Zermelo-Fraenkel-Mengenlehre, die wiederum über ein Axiomensystem aufgebaut wird. Man könnte also sagen, wir betrachten das über ein Teilsystem von Zermelo-Fraenkel oder nehmen irgendein anderes Logiksystem. Die Idee ist jedenfalls: Wir gehen von dieser rein syntaktischen Ebene weg und schauen uns nun die konkreten Modelle unserer Mathematik an.
\end{spoken-clean}

\begin{didactic-insight}[Syntax vs. Semantik]
Der Dozent verwendet verschiedene Metaphern, um den Unterschied zwischen Syntax und Semantik zu verdeutlichen:
\begin{itemize}
    \item \textbf{Musik:} Die Syntax entspricht der Partitur (den geschriebenen Noten), während die Semantik die tatsächliche Musik ist, die erklingt.
    \item \textbf{Kulinarik:} Die Syntax ist wie die Speisekarte, die Semantik ist die tatsächliche Speise, die auf den Tisch serviert wird.
\end{itemize}
In der Logik bedeutet dies: Syntax ist das formale Manipulieren von Zeichenketten nach festen Regeln, während Semantik die Interpretation dieser Zeichen in einer mathematischen Struktur (einem Modell) ist, in der Aussagen wahr oder falsch sein können.
\end{didactic-insight}

\subsection{Syntax und Semantik im Überblick}

\begin{spoken-clean}[00:35:40 - 00:37:30]
Ich glaube, in diesem Bereich werden alle gerne Metaphern verwendet. Im Skript von Lorenz Halbeisen wird es ein bisschen so beschrieben, dass man sich das wie Partitur und Musik vorstellen kann. Wir haben die Partitur, in der alle Noten stehen; auf der semantischen Ebene hingegen erklingt die Musik. Aber die eigentliche Musik ist eben nochmals etwas anderes.

Oder Richard Pink hat es eher kulinarisch verglichen: Die ganze Syntaktik ist wie die Speisekarte, und die Semantik ist wie die Speise, die auf den Tisch serviert wird. Da dürfen Sie auch gerne Ihre eigene Lieblingsmetapher finden.

Grob gesagt: Wir haben die syntaktische Ebene mit Formeln, Beweisen und Zeichen. Und wir haben die semantische Ebene mit Strukturen, Bedeutung sowie Wahrheit oder Nicht-Wahrheit.
\end{spoken-clean}

\subsection{Gegenüberstellung: Syntax vs. Semantik}
\begin{math-stroke}
\begin{center}
\begin{tabular}{l | l}
    \toprule
    \textbf{Syntaktik (Syntax)} & \textbf{Semantik} \\ 
    \midrule
    Terme & Objekte \\
    Formeln & Aussagen \\
    Logische Axiome & Tautologien \\
    Nicht-logische Axiome & Axiomensystem einer Theorie \\
    \bottomrule
\end{tabular}
\end{center}

\begin{explanation-of-steps}
\begin{itemize}
    \item \textbf{Terme vs. Objekte:} Ein Term ist ein syntaktischer Ausdruck (z.\,B. $x+y$), ein Objekt ist das tatsächliche Element in einer Menge.
    \item \textbf{Formeln vs. Aussagen:} Eine Formel ist eine Zeichenkette, eine Aussage ist etwas, das in einer Struktur wahr oder falsch sein kann.
    \item \textbf{Logische Axiome vs. Tautologien:} Logische Axiome sind formal gesetzte Startpunkte; Tautologien sind Formeln, die in \emph{jeder} denkbaren Struktur wahr sind.
    \item \textbf{Nicht-logische Axiome:} Diese definieren eine spezifische mathematische Theorie (z.\,B. Gruppen oder Körper).
\end{itemize}
\end{explanation-of-steps}
\end{math-stroke}

\begin{spoken-clean}[00:37:30 - 00:39:40]
Ich glaube, das Ganze geht auch ein bisschen auf die Sprachphilosophie zurück. Es ist ein bisschen wie bei der Sprache: Man kann rein mit der Sprache hantieren und sagen: \qt{Wenn es regnet, ist die Straße nass}. Damit kann man rein logisch herumspielen und Dinge folgern. In der Mathematik reicht das vielleicht auch. Aber man kann eben auch sagen: Hier ist die Straße, hier ist der Regen, und das bedeutet es, nass zu sein. Dann kann man trotzdem logisch operieren und etwas über die Straße aussagen. Und genau das tun wir hier.

Machen wir das doch präzise --- aber eben im Kontext einer naiven Mengenlehre und eines gewissen Wahrheitsbegriffs, den wir bereits haben. Man könnte philosophisch wahrscheinlich noch viel länger erörtern, was genau dieser Wahrheitsbegriff ist. Wie auch immer, ich denke, es wird alles klar werden.
\end{spoken-clean}

\subsection{\texorpdfstring{$\mathcal{L}$}{L}-Strukturen}

\begin{spoken-clean}[00:39:40 - 00:41:30]
Noch eine Definition. Wir haben eine Signatur $\mathcal{L}$. Wir sagen nun: Eine $\mathcal{L}$-Struktur $M$ besteht aus einer nicht-leeren Menge $A$ und einer Abbildung. Diese Abbildung ordnet jedem Konstantensymbol $c \in \mathcal{L}$ ein Element $c^M \in A$ zu. Wir verwenden hier auch den Begriff des Elementseins einer Menge, aber eben im naiven Sinn.

Jedem $n$-stelligen Relationssymbol $R \in \mathcal{L}$ ordnen wir eine Menge von $n$-Tupeln $R^M$ zu. Ein $n$-Tupel ist dabei eine Teilmenge von $A^n$ [Anmerkung: Ein $n$-Tupel ist ein Element von $A^n$; eine Menge von $n$-Tupeln ist eine Teilmenge von $A^n$]. Wir sagen einfach: All diese $n$-Tupel sind diejenigen, die diese Relation erfüllen. Das ist einfach eine Mengenschreibweise für Relationen, und die verwenden wir hier. Für jedes Relationssymbol erhalten wir also tatsächlich eine Relation auf unserem Bereich $A$.
\end{spoken-clean}

\begin{math-stroke}[Definition der \texorpdfstring{$\mathcal{L}$}{L}-Struktur]
\begin{definition}[$\mathcal{L}$-Struktur]\label[definition]{def:l-struktur}
Sei $\mathcal{L}$ eine Signatur. Eine \newterm{$\mathcal{L}$-Struktur} $M$ besteht aus:
\begin{enumerate}[label=\textbf{(\arabic*)}]
    \item Einer nicht-leeren Menge $A$, genannt der \newterm{Bereich} (oder das \newterm{Universum}) von $M$.
    \item Einer Zuordnung, die:
    \begin{itemize}
        \item jedem Konstantensymbol $c \in \mathcal{L}$ ein Element $c^M \in A$ zuordnet,
        \item jedem $n$-stelligen Relationssymbol $R \in \mathcal{L}$ eine Relation $R^M \subseteq A^n$ zuordnet,
        \item jedem $n$-stelligen Funktionssymbol $F \in \mathcal{L}$ eine Funktion $F^M \colon A^n \to A$ zuordnet.
    \end{itemize}
\end{enumerate}
\end{definition}

\begin{remark}[Relationen als Teilmengen]
Eine $n$-stellige Relation auf einer Menge $A$ ist formal nichts anderes als eine Teilmenge $R \subseteq A^n$.
\begin{example}[Die Kleiner-Relation]
Die strikte Kleiner-Relation $<$ auf $\mathbb{R}$ entspricht der Teilmenge:
\[ R_< = \{ (x, y) \in \mathbb{R}^2 \mid x < y \} \subseteq \mathbb{R}^2 \]
\end{example}
\end{remark}
\end{math-stroke}

\begin{spoken-clean}[00:41:30 - 00:42:00]
Und jedem $n$-stelligen Funktionssymbol $F$ in unserer Signatur $\mathcal{L}$\dots was ordnet man einem Funktionssymbol sinnvollerweise zu?
\end{spoken-clean}

\begin{student-interaction}[Student Answer]
Einen Funktionswert.
\end{student-interaction}

\begin{spoken-clean}[00:42:00 - 00:42:50]
Eine Funktion! Einfach eine Funktion. Es ordnet ihm eine Funktion $F^M \colon A^n \to A$ zu. Und diese Menge $A$ heißt der Bereich von $M$. Genau, das ist eine $\mathcal{L}$-Struktur.

Vielleicht ein kleines Beispiel. Schauen wir uns unsere Signatur an; das soll die Signatur der Gruppentheorie sein. Was war die Signatur der Gruppentheorie?
\end{spoken-clean}

\begin{student-interaction}[Student Answer]
$e$ und mal.
\end{student-interaction}

\begin{spoken-clean}[00:42:50 - 00:43:20]
Genau, wir haben nur das Konstantensymbol $e$ und ein zweistelliges Funktionssymbol $\circ$.
\end{spoken-clean}

\subsection{Beispiel: Struktur der Gruppentheorie}
\begin{math-stroke}
Sei $\mathcal{L} = \mathcal{L}_{GT} = \{e, \circ\}$ die Signatur der Gruppentheorie.
Wir definieren eine $\mathcal{L}$-Struktur $M$ mit Bereich $A = \mathbb{Z}$:
\begin{itemize}
    \item \textbf{Konstante:} $e^M := 0 \in \mathbb{Z}$.
    \item \textbf{Funktion:} $\circ^M \colon \mathbb{Z} \times \mathbb{Z} \to \mathbb{Z}, \quad (a, b) \mapsto a + b$.
\end{itemize}

\begin{explanation-of-steps}
Diese Struktur $M$ interpretiert das neutrale Element $e$ als die Zahl $0$ und die Gruppenoperation $\circ$ als die gewöhnliche Addition ganzer Zahlen. Da $(\mathbb{Z}, 0, +)$ eine Gruppe ist, erfüllt diese Struktur tatsächlich alle Gruppenaxiome. Beachten Sie, dass wir a priori auch $e^M = 100$ und $\circ^M(a, b) = a \cdot b + 3$ hätten wählen können --- dies wäre immer noch eine gültige $\mathcal{L}$-Struktur, auch wenn sie die Gruppenaxiome vermutlich nicht erfüllen würde.
\end{explanation-of-steps}
\end{math-stroke}

\begin{spoken-clean}[00:43:20 - 00:45:10]
Nehmen wir zum Beispiel als Bereich $A$ die ganzen Zahlen $\mathbb{Z}$. Wir definieren $e^M$ als $0$. Dem zweistelligen Funktionssymbol $\circ^M$ ordnen wir die Funktion $\mathbb{Z} \times \mathbb{Z} \to \mathbb{Z}$ zu, die ein Paar $(a, b)$ auf $a + b$ abbildet. Das definiert eine $\mathcal{L}$-Struktur.

Hier ist es besonders schön, weil das tatsächlich eine Gruppe ist, wie wir wissen. Das heißt, diese Struktur erfüllt auch gleich alle Gruppenaxiome. Aber a priori haben wir das überhaupt noch nicht gefordert. Möglicherweise hätte ich auch sagen können, $e^M$ ist $100$ und die Verknüpfung ist irgendwie Multiplikation plus $3$ oder so. Das wäre kein Problem, es muss nur eine Funktion sein.
\end{spoken-clean}

\begin{lecture-break}[Pause]
Der Dozent kündigt eine Pause an. Der Video-Stream zeigt für einige Minuten ein Standbild der Tafel.
\end{lecture-break}

\subsection{Variablenbelegungen und Interpretationen}

\begin{spoken-clean}[00:48:00 - 00:50:35]
Wir machen weiter. Bitte nehmen Sie Ihren Platz ein und beenden Sie Ihre Konversationen. Also einfach nochmals --- Entschuldigung, ich habe es blöderweise schon weggewischt ---: Was ist eine $\mathcal{L}$-Struktur? Eine $\mathcal{L}$-Struktur bedeutet: Wir haben einen Bereich von $M$, das ist eine Menge $A$. Jedem Konstantensymbol ist ein Element im Bereich $A$ zugeordnet. Das heißt, die Konstanten sind jetzt wirklich Elemente. Die Funktionen sind wirklich Funktionen von unserem Bereich (oder dem kartesischen Produkt von $A$ mit sich selbst) nach $A$. Und die Relationssymbole sind wirklich Relationen zwischen den Elementen von $A$.

Jetzt wollen wir noch genauer festlegen, was ein Modell ist. Es kommen jetzt wieder --- ein bisschen wie in der ersten Stunde, als wir die Syntax eingeführt haben --- eine Reihe von Definitionen. Sie sind aber alle recht natürlich; man muss sie sich einfach nochmals in Ruhe anschauen.
\end{spoken-clean}

\begin{spoken-clean}[00:50:35 - 00:53:05]
Eine \newterm{Variablenbelegung} $j$ einer $\mathcal{L}$-Struktur $M$ mit Bereich $A$ ist eine Abbildung, die jeder Variablen $\nu$ ein Objekt $j(\nu) \in A$ zuordnet. Was könnte eine Variablenbelegung sein, rein vom Wort her? Genau, wir wollen jede Variable mit einer Konstanten belegen. Wir sagen also, jede Variable soll für ein Element aus unserem Bereich $A$ stehen. Da gibt es natürlich viele Möglichkeiten.

Es hilft, Folgendes zu schreiben: Wenn $\nu$ eine Variable und $j$ eine Variablenbelegung ist, definieren wir eine neue Variablenbelegung, indem wir $j$ ein wenig abändern. Wir definieren nun $j \frac{a}{\nu}$ für ein $a \in A$. Das Ergebnis ist $a$, falls die betrachtete Variable genau $\nu$ ist, und ansonsten bleibt es einfach bei $j(\nu')$. Das heißt, wir nehmen die Variablenbelegung $j$ und verändern sie nur für die Variable $\nu$. Anstatt $j(\nu)$ nehmen wir nun $a$, aber ansonsten bleibt alles gleich. Wir ändern die Belegung also nur für eine einzige Variable.
\end{spoken-clean}

\begin{math-stroke}[Variablenbelegung]
\begin{definition}[Variablenbelegung]\label[definition]{def:variablenbelegung}
Sei $M$ eine $\mathcal{L}$-Struktur mit Bereich $A$. Eine \newterm{Variablenbelegung} $j$ ist eine Abbildung:
\[ j \colon \text{Var} \to A \]
die jeder Variablen $\nu$ ein Element $j(\nu) \in A$ zuordnet.
\end{definition}

\begin{definition}[Modifizierte Variablenbelegung]\label[definition]{def:mod-variablenbelegung}
Sei $j$ eine Variablenbelegung, $\nu$ eine Variable und $a \in A$. Die \newterm{modifizierte Variablenbelegung} $j \frac{a}{\nu}$ ist definiert durch:
\[
j \frac{a}{\nu}(\nu') := \begin{cases}
a & \text{falls } \nu' \equiv \nu \\
j(\nu') & \text{sonst}
\end{cases}
\]
\end{definition}

\begin{explanation-of-steps}
Die Notation $j \frac{a}{\nu}$ beschreibt eine Belegung, die fast identisch mit $j$ ist, außer dass sie der spezifischen Variablen $\nu$ den Wert $a$ zuweist, unabhängig davon, was $j(\nu)$ ursprünglich war. Dies ist essenziell für die spätere Definition der Wahrheit von Quantorenformeln.
\end{explanation-of-steps}
\end{math-stroke}

\begin{spoken-clean}[00:53:05 - 00:55:40]
Das ist also nichts Kompliziertes oder Gefährliches, nur etwas mühsam. Gut, und jetzt können wir sagen, was eine Interpretation ist. Eine $\mathcal{L}$-Interpretation $I$ ist einfach ein Paar $(M, j)$, wobei $M$ eine $\mathcal{L}$-Struktur und $j$ eine Variablenbelegung ist. Auch hier können Sie wieder versuchen, schöne Metaphern zu finden --- Sprache, Modell und so weiter.

Aber hier sind wir: Wir haben eine $\mathcal{L}$-Struktur $M$ und sagen zusätzlich für jede Variable, welches Element im Bereich von $M$ sie repräsentiert. Und dann können wir weiter definieren: Für eine Interpretation $I = (M, j)$ definieren wir die Interpretation $I \frac{a}{\nu}$ einfach als die Interpretation, die durch $M$ gegeben ist, aber mit der modifizierten Variablenbelegung $j \frac{a}{\nu}$.
\end{spoken-clean}

\subsubsection{Interpretation}
\begin{math-stroke}
\begin{definition}[$\mathcal{L}$-Interpretation]\label[definition]{def:l-interpretation}
Eine \newterm{$\mathcal{L}$-Interpretation} $I$ ist ein geordnetes Paar:
\[ I = (M, j) \]
bestehend aus einer $\mathcal{L}$-Struktur $M$ und einer Variablenbelegung $j$.
\end{definition}

\begin{definition}[Modifizierte Interpretation]\label[definition]{def:mod-interpretation}
Für eine Interpretation $I = (M, j)$ definieren wir:
\[ I \frac{a}{\nu} := (M, j \frac{a}{\nu}) \]
\end{definition}
\end{math-stroke}

\subsection{Interpretation von Termen}

\begin{spoken-clean}[00:55:40 - 00:58:20]
Gut, wir haben Struktur, Variablenbelegung und Interpretationen. Als Nächstes wollen wir einen Term interpretieren. Was war noch mal ein Term? Einen Term haben wir aus Konstanten und Variablen konstruiert, auf die wir Funktionen angewendet haben. Das waren Terme --- und dann wieder Funktionen auf Funktionen und so weiter.

Damit können wir nun mit einer $\mathcal{L}$-Interpretation auch jeden Term interpretieren. Das ist die Interpretation eines Terms. Sei also $I = (M, j)$ eine $\mathcal{L}$-Interpretation. Wir ordnen nun jedem Term $\tau$ einen Wert $I(\tau) \in A$ zu.
\end{spoken-clean}

\begin{math-stroke}[Interpretation eines Terms]
Sei $I = (M, j)$ eine $\mathcal{L}$-Interpretation mit Bereich $A$. Wir ordnen jedem $\mathcal{L}$-Term $\tau$ induktiv einen Wert $I(\tau) \in A$ zu:
\begin{enumerate}[label=\textbf{(\arabic*)}]
    \item \textbf{Variablen:} Falls $\tau$ eine Variable $\nu$ ist, dann gilt:
    \[ I(\nu) := j(\nu) \]
    \item \textbf{Konstanten:} Falls $\tau$ ein Konstantensymbol $c$ ist, dann gilt:
    \[ I(c) := c^M \]
    \item \textbf{Funktionen:} Falls $\tau$ von der Form $F \tau_1 \dots \tau_n$ ist, wobei $F$ ein $n$-stelliges Funktionssymbol und $\tau_1, \dots, \tau_n$ Terme sind, dann gilt:
    \[ I(F \tau_1 \dots \tau_n) := F^M(I(\tau_1), \dots, I(\tau_n)) \]
\end{enumerate}

\begin{explanation-of-steps}
Die Interpretation eines Terms erfolgt induktiv über den Aufbau des Terms:
\begin{itemize}
    \item Variablen erhalten ihren Wert direkt aus der Variablenbelegung $j$.
    \item Konstanten erhalten ihren festen Wert $c^M$ aus der Struktur $M$.
    \item Komplexe Terme werden berechnet, indem man zuerst die inneren Terme $\tau_i$ interpretiert und dann die Funktion $F^M$ der Struktur auf diese Werte anwendet.
\end{itemize}
\end{explanation-of-steps}
\end{math-stroke}

\begin{spoken-clean}[00:58:20 - 01:01:05]
Wie machen wir das? Induktiv natürlich! Erstens: Falls $\tau$ eine Variable $\nu$ ist, dann ist $I(\nu)$ einfach $j(\nu)$. Das ist genau das, was die Variablenbelegung macht. Zweitens: Falls $\tau$ ein Konstantensymbol $c$ ist, dann ist $I(c)$ einfach $c^M$. Das ist genau das, was die Struktur macht.

Und drittens: Falls $\tau$ von der Form $F \tau_1 \dots \tau_n$ ist --- wir haben also ein $n$-stelliges Funktionssymbol und darauf Terme angewendet ---, was machen wir dann? Wir nehmen die Funktion $F^M$, die uns die Struktur gibt, und wenden sie auf die Interpretationen der einzelnen Terme an. Also $F^M$ angewendet auf $I(\tau_1)$ bis $I(\tau_n)$.

Das ist eine rekursive Definition. Wir fangen bei den einfachsten Termen an --- Variablen und Konstanten --- und bauen uns dann zu den komplizierteren Termen hoch. Das ist die Interpretation eines Terms.
\end{spoken-clean}

\begin{spoken-clean}[01:01:05 - 01:04:37]
Gut, jetzt haben wir Terme interpretiert. Jetzt kommt der nächste Schritt: Wir wollen Formeln interpretieren. Wir wollen sagen können, wann eine Formel in einer Interpretation wahr ist. Das ist die berühmte Tarski-Definition der Wahrheit. Alfred Tarski war ein polnischer Logiker, der das in den 1930er Jahren formalisiert hat.

Wir schreiben dafür $I \models \varphi$ und lesen das als: \qt{Die Interpretation $I$ ist ein Modell für $\varphi$} oder \qt{$\varphi$ ist wahr in $I$}. Auch das machen wir wieder induktiv über den Aufbau der Formeln. Wir fangen an mit den atomaren Formeln. Das waren Gleichungen zwischen Termen oder Relationen zwischen Termen.

Also: $I \models \tau_1 = \tau_2$ gilt genau dann, wenn die Interpretationen der Terme gleich sind, also wenn $I(\tau_1) = I(\tau_2)$ als Elemente in unserem Bereich $A$ gilt. Und für ein Relationssymbol: $I \models R \tau_1 \dots \tau_n$ gilt genau dann, wenn das Tupel der Interpretationen $(I(\tau_1), \dots, I(\tau_n))$ in der Relation $R^M$ liegt. Das ist der Anfang. Und dann kommen die logischen Verknüpfungen und die Quantoren.
\end{spoken-clean}

\section{Wahrheitsbegriff in der Modelltheorie}

\begin{spoken-clean}[01:04:37 - 01:06:57]
Gut, nach den Termen haben wir Formeln definiert. Was wir jetzt tun wollen: Wir wollen sagen, wann eine Formel in einer gewissen Interpretation wahr ist oder nicht. Das können wir jetzt auch wieder rekursiv definieren.

Als kurze Erinnerung: Wie sind Formeln entstanden? Wir haben Terme genommen und sie durch die Regeln $F_0$ bis $F_4$ zusammengesetzt. Wir haben gesagt: Wenn ein Term gleich einem anderen Term ist, ist das eine Formel. Oder wenn wir ein Relationssymbol hatten, konnten wir Terme einsetzen. Danach konnten wir aus bestehenden Formeln neue bauen, indem wir logische Quantoren und Junktoren verwendet haben.

Nochmals zur Erinnerung: Eine Formel $\varphi$ ist aus den Regeln $F_0$ bis $F_4$ entstanden. Das heißt, eine Formel hat die Form $\tau_1 = \tau_2$, oder sie ist eine Relation zwischen Termen. Wenn $\psi$ eine Formel ist, kann es auch $\neg \psi$ sein, oder $\psi_1 \land \psi_2, \psi_1 \lor \psi_2, \psi_1 \to \psi_2$, oder eben mit den Quantoren $\exists \nu \psi$ oder $\forall \nu \psi$.
\end{spoken-clean}

\begin{math-stroke}[Wahrheit einer Formel]
\begin{definition}[Wahrheitsdefinition nach Tarski]\label[definition]{def:tarski-wahrheit}
Für eine Formel $\varphi$ definieren wir die Relation \qt{$I \models \varphi$} (gelesen: \qt{$\varphi$ gilt in $I$} oder \qt{$\varphi$ ist wahr bezüglich $I$}) rekursiv wie folgt:
\begin{enumerate}[label=\textbf{(W\arabic*)}, start=0, itemsep=0.8ex, labelsep=0.8em, labelindent=2em, leftmargin=*]
    \item \textbf{Gleichheit:} $I \models \tau_1 = \tau_2 \iff I(\tau_1) = I(\tau_2)$ (als Objekte in $A$).
    \item \textbf{Relationen:} $I \models R \tau_1 \dots \tau_n \iff (I(\tau_1), \dots, I(\tau_n)) \in R^M$.
    \item \textbf{Negation:} $I \models \neg \psi \iff \text{nicht } I \models \psi$.
    \item \textbf{Konjunktion:} $I \models \psi_1 \land \psi_2 \iff I \models \psi_1 \text{ und } I \models \psi_2$.
    \item \textbf{Disjunktion:} $I \models \psi_1 \lor \psi_2 \iff I \models \psi_1 \text{ oder } I \models \psi_2$.
    \item \textbf{Implikation:} $I \models \psi_1 \to \psi_2 \iff (\text{falls } I \models \psi_1, \text{ dann } I \models \psi_2)$.
    \item \textbf{Existenzquantor:} $I \models \exists \nu \psi \iff \exists a \in A \colon I \frac{a}{\nu} \models \psi$.
    \item \textbf{Allquantor:} $I \models \forall \nu \psi \iff \forall a \in A \colon I \frac{a}{\nu} \models \psi$.
\end{enumerate}
\end{definition}
\end{math-stroke}

\begin{spoken-clean}[01:06:57 - 01:09:17]
Das ist jetzt eine längliche Definition. Wir wollen all das in die semantische Ebene übersetzen. Das tun wir, indem wir genau das verwenden, was diese Schriftzeichen bedeuten sollen.

Wir sagen zuerst einmal, dass $\tau_1 = \tau_2$ wahr ist bezüglich $I$, falls --- per Definition --- die Interpretation von $\tau_1$ (ein Element in $A$) dasselbe Objekt ist wie die Interpretation von $\tau_2$. Das ergibt absolut Sinn. $\tau_1 = \tau_2$ ist hier erst einmal nur ein Symbolzeichen. Wenn wir aber eine Interpretation haben, werden daraus zwei Elemente, und wir können schauen: Sind diese Elemente dieselben Objekte? Falls ja, dann schreiben wir das so hin.

Dann schreiben wir: In $I$ gilt die Relation $R(\tau_1, \dots, \tau_n)$ genau dann, wenn die entsprechenden Elemente diese Relation erfüllen. Das heißt, wenn dieses Tupel in unserer Relation $R^M$ enthalten ist --- also in dieser Teilmenge von $A^n$, die durch unsere Struktur $M$ gegeben ist.
\end{spoken-clean}

\begin{spoken-clean}[01:09:17 - 01:11:47]
Das andere machen wir jetzt wieder rekursiv. Wir sagen, dass $\neg \psi$ wahr ist in $I$ genau dann, wenn $\psi$ in $I$ nicht wahr ist. Wir können also prüfen: Wenn das nicht gilt, dann schreiben wir das so. Hier haben wir einen naiven Wahrheitsbegriff, um zu sagen, ob zwei Elemente in einer Menge gleich sind oder nicht. Und jetzt gehen wir weiter durch, indem wir genau das tun, was wir vielleicht schon immer klammheimlich interpretiert haben.

Wir sagen jetzt: $\psi_1 \land \psi_2$ gilt in $I$ genau dann, wenn $\psi_1$ in $I$ gilt und $\psi_2$ in $I$ gilt. Für das Oder: $\psi_1 \lor \psi_2$ gilt genau dann, wenn $\psi_1$ gilt oder $\psi_2$ gilt. Und für die Implikation: $\psi_1 \to \psi_2$ gilt genau dann, wenn falls $\psi_1$ in $I$ gilt, dann auch $\psi_2$ in $I$ gilt.
\end{spoken-clean}

\begin{spoken-clean}[01:11:47 - 01:15:07]
Dann brauchen wir noch die Quantoren. Da benötigen wir unsere Notation: Wir sagen, $\exists \nu \psi$ gilt in $I$ --- das entspricht nun: Es existiert ein $a$ in $A$, so dass\dots jetzt wird es ein bisschen technischer\dots wenn wir in dieser Interpretation die Variable $\nu$ durch $a$ ersetzen, dann gilt $\psi$ in dieser modifizierten Interpretation.

Und wir sagen, dass $\forall \nu \psi$ genau dann gelten soll, wenn für alle $a$ in $A$ gilt: Wenn wir diese andere Variablenbelegung nehmen, bei der wir $\nu$ durch $a$ ersetzen, soll immer noch $\psi$ gelten.

Das ist diese Definition. Es ist genau das, was man denkt. Links steht unser formales System, unsere formale Sprache. Das auf der rechten Seite sind jetzt gewissermaßen metamathematische Bedingungen. Sie liegen außerhalb des formalen Systems, das wir aufgebaut haben. Im Skript und im Buch machen sie es so, dass sie diese Bedingungen in Großbuchstaben schreiben, um zu signalisieren: Das ist Metamathematik, bei der man einen Schritt zurücktritt. Aber einfach gesagt: Hier haben wir einen anderen Begriff. Vorher hatten wir reine syntaktische Zeichen, und hier haben wir jetzt eine Interpretation.
\end{spoken-clean}

\begin{didactic-insight}[Objektsprache vs. Metasprache]
Die Tarski-Definition schlägt die Brücke zwischen der formalen Sprache (Objektsprache) und der Sprache, in der wir über Mathematik reden (Metasprache). Wenn wir schreiben $I \models \psi_1 \land \psi_2 \iff I \models \psi_1 \text{ und } I \models \psi_2$, dann ist das \qt{$\land$} ein Symbol der Logik, während das \qt{und} ein Wort unserer natürlichen Sprache ist. Tarski formalisierte dies, um Paradoxien zu vermeiden und den Wahrheitsbegriff mathematisch präzise zu fassen.
\end{didactic-insight}

\subsection{Modelle}

\begin{spoken-clean}[01:15:07 - 01:17:32]
Wir schreiben $I \not\models \varphi$, falls $\varphi$ in der Interpretation $I$ nicht gilt. Es ist immer so, dass in $I$ entweder $\varphi$ gilt oder --- per Definition, wie wir Nicht-$\varphi$ festgelegt haben --- wenn $\varphi$ nicht gilt, dann muss $\neg \varphi$ gelten. Es ist ein exklusives Oder, das heißt, es gilt nie beides.

Manche Leute verwenden die Wendung \qt{entweder oder}, um ein exklusives Oder auszudrücken, aber das ist oft verwirrend. Das sollte man besser nicht tun. Ja? Bitte?
\end{spoken-clean}

\begin{student-interaction}[Studentenfrage]
Ist das nicht dasselbe wie die Definition der Negation vorhin?
\end{student-interaction}

\begin{spoken-clean}[continued]
Nein, es ist genau das Gleiche, aber es ist hier Teil der rekursiven Definition. Es ist eigentlich eine Übersetzung von dem einen in das andere. Man verwendet diese Schreibweise manchmal. Aber das geschieht innerhalb der Formel --- wenn Sie eine Formel haben, gehen Sie diese Schritt für Schritt durch. Wenn Sie zum Beispiel haben \qt{für alle $x$ gilt: nicht $\varphi$ impliziert $\psi$}, dann kommt das irgendwann an die Reihe.

Genau, also Achtung mit \qt{entweder oder}. Ich würde das im mathematischen Bereich eher vermeiden und, wenn man ein exklusives Oder meint, das ausdrücklich so sagen.
\end{spoken-clean}

\begin{spoken-clean}[01:17:32 - 01:20:07]
Gut, jetzt haben wir endlich alles beisammen, um sagen zu können, was ein Modell ist. Sei $\mathcal{L}$ eine Signatur, $\varphi$ eine $\mathcal{L}$-Formel und $M$ eine $\mathcal{L}$-Struktur. Dann sagen wir, $M$ ist ein \newterm{Modell} von $\varphi$ --- wir schreiben $M \models \varphi$ ---, falls für jede Variablenbelegung $j$ gilt, dass die Interpretation $(M, j) \models \varphi$ erfüllt.

Das heißt: Egal wie Sie die Variablen belegen, es soll immer wahr sein. Wenn eine Formel freie Variablen hat, dann muss sie für jede mögliche Belegung dieser Variablen im Bereich $A$ wahr sein, damit die Struktur ein Modell ist.

Dasselbe kann man natürlich auch für eine Menge von $\mathcal{L}$-Formeln machen. Sei $\Phi$ eine Menge von $\mathcal{L}$-Formeln. Dann ist $M$ ein Modell von $\Phi$, falls für jede Formel $\varphi \in \Phi$ gilt, dass $M$ ein Modell von $\varphi$ ist.
\end{spoken-clean}

\begin{math-stroke}[Definition des Modells]
\begin{definition}[Modell]\label[definition]{def:modell}
Sei $\mathcal{L}$ eine Signatur, $\varphi$ eine $\mathcal{L}$-Formel und $M$ eine $\mathcal{L}$-Struktur.
\begin{enumerate}[label=\textbf{(\arabic*)}]
    \item $M$ heißt ein \newterm{Modell} von $\varphi$ (geschrieben $M \models \varphi$), falls für \emph{jede} Variablenbelegung $j$ gilt:
    \[ (M, j) \models \varphi \]
    \item Sei $\Phi$ eine Menge von $\mathcal{L}$-Formeln. $M$ heißt ein \newterm{Modell} von $\Phi$ (geschrieben $M \models \Phi$), falls für jede Formel $\varphi \in \Phi$ gilt:
    \[ M \models \varphi \]
\end{enumerate}
\end{definition}
\end{math-stroke}

\section{Beispiele für Modelle}

\begin{spoken-clean}[01:20:07 - 01:22:55]
Machen wir noch Beispiele. Wir hatten vorher kurz die Gruppen\-theorie-\hspace{0pt}Struktur angeschaut, die durch die ganzen Zahlen $\mathbb{Z}$ gegeben ist. Wir hatten diese $\mathcal{L}$-\hspace{0pt}Gruppen\-theorie-\hspace{0pt}Struktur: Der Bereich $A$ sind die ganzen Zahlen $\mathbb{Z}$. Wir haben diese Struktur $M$ genannt. Dann haben wir gesagt, dass wir das Konstanten\-symbol $e^M$ der $0$ zuordnen. Und wir sagen, diese Verknüpfung $\circ^M$ soll $a + b$ sein.

Dann ist das ein Modell, wenn wir für $\Phi$ die Axiome der Gruppentheorie nehmen. Das kann man sagen, weil all diese Formeln in $\Phi$ --- also die Gruppentheorie --- erfüllt werden. Egal mit welcher Variablenbelegung: Diese Verknüpfung hier erfüllt alle Axiome der Gruppentheorie. Man kann alle Axiome durchgehen und sehen, dass das stimmt, egal für welche Variablenbelegung. Also ist das ein Modell der Gruppentheorie.
\end{spoken-clean}

\begin{math-stroke}[Beispiel: Modell der Gruppentheorie]
\begin{example}[Die additive Gruppe der ganzen Zahlen]\label[example]{ex:gruppe-z}
Sei $\mathcal{L}_{GT} = \{e, \circ\}$ die Signatur der Gruppentheorie und $\Phi_{GT}$ die Menge der Gruppenaxiome.
Wir betrachten die Struktur $M = (\mathbb{Z}, 0, +)$.
Es gilt:
\[ M \models \Phi_{GT} \]
\end{example}

\begin{explanation-of-steps}
Da $(\mathbb{Z}, +)$ mit dem neutralen Element $0$ eine abelsche Gruppe bildet, sind alle Gruppenaxiome (Assoziativität, Existenz des Neutralen, Existenz der Inversen) für jede Belegung der Variablen mit ganzen Zahlen erfüllt. Somit ist diese Struktur ein Modell für die Theorie der Gruppen.
\end{explanation-of-steps}
\end{math-stroke}

\begin{spoken-clean}[01:22:55 - 01:25:32]
Wir können auch einfachere Beispiele anschauen. Machen wir da irgendetwas... Man kann das ein bisschen selbst ausprobieren. Sagen wir, die Signatur $\mathcal{L}$ sei einfach $\{c, F\}$. Wir legen fest: $c$ ist ein Konstantensymbol und $F$ soll ein einstelliges Funktionssymbol sein.

Nun nehmen wir eine Theorie $\Phi$, die aus zwei Formeln besteht. Die eine Formel sei $\varphi_1$: Für alle $x$ gilt $x = c$ oder $x = F(c)$. Und $\varphi_2$ sei: Es existiert ein $x$, so dass $x \neq c$ ist. Das ist jetzt einfach eine etwas zufällige Theorie ohne tiefen Hintergrund, aber man kann das so formulieren.

Die Frage ist jetzt: Was haben wir dafür für Modelle? Wir nehmen zwei verschiedene $\mathcal{L}$-Strukturen, beide mit demselben Bereich. Nehmen wir als Bereich $A$ einfach die Menge bestehend aus zwei Elementen, $0$ und $1$. Wir definieren nun auf diesem Bereich zwei $\mathcal{L}$-Strukturen $M_1$ und $M_2$.
\end{spoken-clean}

\begin{math-stroke}[Beispiel: Eine kleine Theorie]
\textbf{Signatur:} $\mathcal{L} = \{c, F\}$ ($c$ Konstante, $F$ 1-stelliges Funktionssymbol).
\textbf{Theorie:} $\Phi = \{\varphi_1, \varphi_2\}$ mit:
\begin{align*}
    \varphi_1 &\equiv \forall x \, (x = c \lor x = F(c)) \\
    \varphi_2 &\equiv \exists x \, (x \neq c)
\end{align*}

\textbf{Strukturen:} Bereich $A = \{0, 1\}$.
\begin{itemize}
    \item \textbf{Struktur $M_1$:} $c^{M_1} := 0$, $F^{M_1}(0) := 1$, $F^{M_1}(1) := 0$.
    \item \textbf{Struktur $M_2$:} $c^{M_2} := 0$, $F^{M_2}(0) := 0$, $F^{M_2}(1) := 1$.
\end{itemize}
\end{math-stroke}

\begin{spoken-clean}[01:25:32 - 01:27:03]
Und zwar folgendermaßen: Die Konstante in $M_1$ definieren wir als $0$. Dann müssen wir die Funktion in der ersten Struktur definieren. Da sagen wir einfach: Die Funktion ist einstellig; sie bildet $0$ auf $1$ ab und $1$ auf $0$. Das ist eine gültige $\mathcal{L}$-Struktur.

Die zweite Struktur $M_2$: Da sei die Konstante wieder $0$. Die Funktion $F^{M_2}$ definieren wir etwas anders: Sagen wir, $0$ soll auf $0$ abgebildet werden und $1$ soll auf $1$ abgebildet werden.

Jetzt müssen wir schauen: Was gilt hier? Wir haben auf jeden Fall, dass $M_1$ ein Modell von $\varphi_2$ ist, und $M_2$ ist ebenfalls ein Modell von $\varphi_2$. Denn $\varphi_2$ sagt einfach: Es existiert ein $x$, so dass $x \neq c$ ist. $c$ ist hier einfach $0$, und in beiden Modellen existiert ein Element, das nicht $0$ ist, nämlich die $1$. Das ist gut.

Aber für die zweite Formel sehen wir, dass $M_1$ ein Modell von $\varphi_1$ ist. Wir haben: Jedes Element $x$ ist entweder gleich der Konstanten oder gleich dem Funktionswert der Konstanten. Das passt: Entweder ist es $0$ oder es ist $1$. Aber hier in $M_2$ ist der Funktionswert der Konstanten auch wieder $0$. Das heißt, das Element $1$ erfüllt das nicht. In $M_2$ gilt $\varphi_1$ also nicht.

So viel zur Modelltheorie. Vielleicht noch eine Bemerkung: Modelltheorie ist wirklich ein eigenes Gebiet der Mathematik; das hier ist erst der ganz Anfang. Es sind viele Bücher dazu geschrieben worden, und es ist ein aktives Forschungsgebiet. Man verwendet sie effektiv auch, um mathematische Sätze zu beweisen --- es gibt Bereiche der Zahlentheorie, die viel Modelltheorie verwenden, um neue Sätze zu beweisen. Okay. Vielen Dank fürs Kommen, und wir sehen uns nächste Woche wieder.
\end{spoken-clean}

\begin{math-stroke}[Analyse der Modelle]
\textbf{Prüfung von $\varphi_2 \equiv \exists x \, (x \neq c)$:}
\begin{itemize}
    \item In $M_1$: $c^{M_1} = 0$. Für $x=1$ gilt $1 \neq 0$. Also $M_1 \models \varphi_2$.
    \item In $M_2$: $c^{M_2} = 0$. Für $x=1$ gilt $1 \neq 0$. Also $M_2 \models \varphi_2$.
\end{itemize}

\textbf{Prüfung von $\varphi_1 \equiv \forall x \, (x = c \lor x = F(c))$:}
\begin{itemize}
    \item In $M_1$: $c^{M_1} = 0$ und $F^{M_1}(c^{M_1}) = F^{M_1}(0) = 1$. Da $A = \{0, 1\}$, ist jedes Element entweder $0$ oder $1$. Also $M_1 \models \varphi_1$.
    \item In $M_2$: $c^{M_2} = 0$ und $F^{M_2}(c^{M_2}) = F^{M_2}(0) = 0$. Die Bedingung lautet also $x = 0 \lor x = 0$. Für $x=1$ ist dies falsch. Also $M_2 \not\models \varphi_1$.
\end{itemize}

\textbf{Fazit:}
$M_1$ ist ein Modell der Theorie $\Phi$ ($M_1 \models \Phi$), aber $M_2$ ist kein Modell der Theorie $\Phi$ ($M_2 \not\models \Phi$).
\end{math-stroke}

\begin{nice-box}[Zusammenfassung der Schlüsselkonzepte]
In dieser Lektion haben wir die Grundlagen der Peano-Arithmetik und den Einstieg in die Modelltheorie behandelt:
\begin{itemize}
    \item \textbf{Peano-Arithmetik:} Ein Axiomensystem für die natürlichen Zahlen basierend auf der Nachfolgerfunktion $S$. Zentral ist das \textbf{Induktionsschema}, das für jede Formel $\varphi$ ein Axiom bereitstellt.
    \item \textbf{Semi-formale Beweise:} Eine praxisnahe Form der Beweisführung, die logische Standardumformungen unterdrückt und sich auf theoriespezifische Axiome konzentriert.
    \item \textbf{Modelltheorie:} Der Übergang von Syntax (Formeln) zu Semantik (Wahrheit). Eine \textbf{$\mathcal{L}$-Struktur} interpretiert die Symbole einer Signatur in einem Bereich $A$.
    \item \textbf{Tarski-Wahrheit:} Eine rekursive Definition, die festlegt, wann eine Formel unter einer Interpretation wahr ist.
    \item \textbf{Modell:} Eine Struktur $M$ ist ein Modell einer Formel $\varphi$, wenn diese für jede Variablenbelegung in $M$ wahr ist.
\end{itemize}
\end{nice-box}